% Ban thu: chi giu cac cau co trong CD2

%@ id=TT25.12.003 | trang=1 | so=2 | de=II | vitri=CD2.4,C3.9.5 | muc=8 | tt=0
\begin{bt}
Cho hai biểu thức: $A=\dfrac{\sqrt{x}+1}{\sqrt{x}-2}$ và $B=\dfrac{\sqrt{x}+2}{\sqrt{x}-3}+\dfrac{\sqrt{x}-8}{(\sqrt{x}-2)(\sqrt{x}-3)}$ với $x\ge 0;\ x\ne 4;\ x\ne 9$.
\begin{cacy}
\item Tính giá trị của biểu thức $A$ khi $x=25$.
\item Rút gọn biểu thức $B$.
\item Tìm tất cả các giá trị nguyên của $x$ để $B<A$.
\end{cacy}
\loigiai{
\begin{cacy}
\item Thay $x=25$ (TMĐK) vào $A$, ta có:
\[
A=\dfrac{\sqrt{25}+1}{\sqrt{25}-2}=\dfrac{6}{3}=2.
\]
Vậy $x=25$ thì $A=2$.

\item Ta có:
\[
B=\dfrac{(\sqrt{x}+2)(\sqrt{x}-2)+\sqrt{x}-8}{(\sqrt{x}-2)(\sqrt{x}-3)}
=\dfrac{x-4+\sqrt{x}-8}{(\sqrt{x}-2)(\sqrt{x}-3)}
=\dfrac{x+4\sqrt{x}-3\sqrt{x}-12}{(\sqrt{x}-2)(\sqrt{x}-3)}
=\dfrac{(\sqrt{x}+4)(\sqrt{x}-3)}{(\sqrt{x}-2)(\sqrt{x}-3)}
=\dfrac{\sqrt{x}+4}{\sqrt{x}-2}.
\]
Vậy, $x\ge 0;\ x\ne 4;\ x\ne 9$ thì $B=\dfrac{\sqrt{x}+4}{\sqrt{x}-2}$.

\item Ta có:
\[
B<A
\]
\[
\dfrac{\sqrt{x}+4}{\sqrt{x}-2}<\dfrac{\sqrt{x}+1}{\sqrt{x}-2}
\]
\[
\dfrac{\sqrt{x}+4}{\sqrt{x}-2}-\dfrac{\sqrt{x}+1}{\sqrt{x}-2}<0
\]
\[
\dfrac{3}{\sqrt{x}-2}<0
\]
\[
\sqrt{x}<2
\]
\[
x<4.
\]
Kết hợp điều kiện, suy ra $0\le x<4$. Mà $x\in\mathbb{Z}$ nên $x\in\{0;1;2;3\}$.

Vậy $x\in\{0;1;2;3\}$ thì $B<A$.
\end{cacy}
}
\end{bt}
